\section{Discussion}
\label{sec:discussion}

The compact selection rule has a simple physical origin. Local Higgs sectors
carry electric charges or flavor moment maps. A compact Calabi--Yau does not
leave all of the corresponding flavor symmetries global: its $C_3$ vectors
gauge diagonal combinations determined by divisor--curve pairings. Their
moment-map equations cut out a correlated subspace of the product of local
deformation coordinates. The same divisor--curve map appears in the
local-to-global deformation sequence, so the five-dimensional vacuum equation
and the geometric smoothing relation have the same kernel.

For nodes this reduces to the familiar conifold transition. The substantive
extension is that an interacting del Pezzo sector enters the relation on the
same footing as a wrapped-M2 hypermultiplet, through a current-multiplet moment
map rather than an elementary field bilinear. The $X_9$ relation
\begin{equation}
 (u_1,u_2,\beta)=t(-1,-1,1)
\end{equation}
is the simplest realization: two conifold condensates and an $E_2$ Higgs
operator participate in one primitive compact transition. Neither part exists
as a global transition without the other.

The $E_3$ sector shows why a branchwise formulation is necessary. Its local
Higgs branch is not irreducible. The components
$\cH_{\mathfrak{sl}_3}$ and $\cH_{\mathfrak{sl}_2}$ couple to different
compact cocharacters, have different dimensions, and can change the rank of
the global gauging matrix. Replacing the local theory by
$\dim T^1=3$ would erase precisely the field-theoretic information responsible
for the profile dependence of $X^{\circ}$.

\subsection{Relation to generalized-symmetry gluing}
\label{subsec:generalized-symmetry-relation}

Compact M-theory models also correlate the generalized symmetries and global
forms of several local sectors. There, Mayer--Vietoris maps determine which
defects and higher-form symmetries glue consistently
\cite{Cvetic:2023plv}. The analogy with the present result is structural. In
both cases, independent local data are replaced by the kernel or cokernel of a
map supplied by the complement of the singular neighborhoods in the compact
threefold.

The constrained objects are different. Generalized-symmetry calculations
track defect groups, polarizations and higher-group structures. Our matrix
acts on ordinary Higgs moment maps and complex deformation coordinates. One
result does not follow from the other. Together they indicate that the global
topology of a compactification controls more than the spectrum of isolated
local theories: it also controls how their moduli and symmetry data can be
assembled.

\subsection{Comparison with the compact del Pezzo--conifold precursor}
\label{subsec:malyshev-comparison}

Malyshev established that a root curve in a compact del Pezzo surface can be
homologous to curves at spatially separated conifold points, and that this
relation reduces the number of del Pezzo deformations
\cite{Malyshev:2007zz}. That calculation proves the geometric phenomenon in a
particular quintic. It does not derive the relation as a five-dimensional
gauging or formulate a theorem for interacting Higgs branches.

The branchwise charge-kernel theorem is therefore a strict extension, both in
scope and in physical content. It applies uniformly to the admissible toric
class and identifies the relation matrix with electric and flavor gaugings.
It retains the reduced and embedded $E_2$ schemes and both $E_3$ components,
determines the exact image of the moment-map locus, and combines it with the
mixed smoothing criterion of \cite{Wuebben:2026smoothing}. The resulting
existence test uses the local discriminants as well as the charge kernel.
The positive and negative examples include
an independent protected-spectrum check. The antecedent is the mixed homology
relation; the new result is its branchwise five-dimensional kernel theorem.

\subsection{Open problems}
\label{subsec:open-problems}

The exact smoothing criterion and selected-base integrability resolve the
geometric existence question under the stated hypotheses, including the
two-dimensional $E_3$ component. They also distinguish abstract smoothings
from a chosen ambient construction. For example, the threefold $X_{19}$
of \cite[Corollary~9.3]{Wuebben:2026smoothing}, with fourteen nodes and
one degree-seven cone point, has a smooth full deformation base of
dimension thirty and admits a smoothing, although every single-degree
ambient family in the specified class fails to smooth it. Failure of such
an ambient construction is therefore not a compact gauging obstruction.

The remaining questions concern the following distinctions.

\begin{enumerate}
 \item \emph{Unrestricted deformation schemes.} Smoothness of a selected
 base does not determine how distinct profiles meet, the full nilpotent
 structure at $E_2$ points, or selected bases with zero $E_3$ projection.
 These are outside the smoothness assertion used here, not gaps in the
 exact criterion for existence of a smoothing.

 \item \emph{Controlled rigid limits.} For $X_9$, the shared flavor gauging
 is lost in every gravity-decoupling limit with bounded local K\"ahler
 scales \cite[Section~4]{Wuebben:2026quantum}. Other compact models,
 including $X^{\circ}$, require their own kinetic and K\"ahler-cone
 analysis; the compact existence criterion alone does not decide which
 gaugings survive a noncompact limit.

 \item \emph{Integral global form.} The examples determine primitive Smith
 data in their toric divisor lattices. Establishing index one in full integral
 cohomology is separate from the rational smoothing criterion. It is
 established for $X_9$ in \cite[Section~3 and Appendix~B]{Wuebben:2026quantum};
 other models require an integral analysis before discrete remnants or
 line operators can be inferred.

 \item \emph{The compact scalar geometry.} Equation
 \eqref{eq:physical-geometric-identification} identifies a protected
 holomorphic projection. Constructing the full coupled
 quaternionic-K\"ahler space and its chiral-ring analogue would go beyond the
 rigid zero-level description.

 \item \emph{Broader sectors.} The geometric theorem already extends beyond
 toric hypersurfaces to projective threefolds with its exact analytic
 germs, including degree-five and $\mathbb P^1\times\mathbb P^1$ cones.
 Extending the local physical dictionaries to those sectors, and then to
 non-isolated singularities and higher-rank theories, requires additional
 operator and charge data.
\end{enumerate}

The mixed-transition mechanism already follows from the results above. Higgs
directions that are independent in isolated local theories need not remain
independent in a compact M-theory vacuum. Dynamical five-dimensional vectors
select their allowed combinations, and a global transition can require
cooperation between free conifold states and an interacting SCFT sector. This
is the compact generalization of the conifold charge relation.
